\documentclass[letterpaper,12pt,oneside]{article}
\usepackage[top=1in, left=1.25in, right=1.25in, bottom=1in]{geometry}

\usepackage[T1]{fontenc}
\usepackage[utf8]{inputenc}
\usepackage[spanish,es-nodecimaldot,es-tabla]{babel}
\usepackage{caption, subcaption}
\usepackage{graphicx}
\usepackage{array}
\usepackage{tikz}
\usepackage{setspace}

\graphicspath{{./figs/}}

\begin{document}

\begin{titlepage}
\setlength{\parindent}{0pt}
\setlength{\parskip}{0pt}

\begin{center}
\vfill

\includegraphics[width=4.5cm]{UNAM_LOGO2.png}

\vspace{0.5cm}

{\Huge\textbf{Universidad Nacional Autónoma de México}}

\vfill

{\large Ingeniería en ...}

\vfill

{\large Estructura de Datos y Algoritmos I}

\vspace{1cm}

{\huge FORMATO DE REPORTE}

\vspace{1cm}

ALUMNO: \\
\large 323197294

\bigskip

Grupo: \\
\#

Semestre: \\
2026-II

\vfill

México, CDMX. Febrero 2026

\end{center}
\end{titlepage}

\tableofcontents
\clearpage

\section{Introducción}

Este apartado debe abordar los siguientes puntos:

\begin{itemize}
    \item \textbf{Planteamiento del problema.}
    \item \textbf{Motivación.}
    \item \textbf{Objetivos.}
\end{itemize}

\section{Marco Teórico}

Un algoritmo se define como un conjunto de instrucciones ordenadas y definidas para resolver un problema a través de secuencias lógicas, es como una receta en la cual se deben seguir ciertos pasos e instrucciones para poder preparar lo que se desea. Este tiene gran importancia ya que es la base lógica para poder resolver problemas y poder llevar a cabo tareas específicas \cite{algoritmo}.

Los arreglos son estructuras de datos que almacenan un conjunto de elementos de un mismo tipo en ubicaciones de memoria contiguas \cite{arreglos}.

Un tipo de dato abstracto o TDA es un modelo matemático que está compuesto por una serie de operaciones que están definidas sobre un conjunto de datos, se caracteriza por tener un tamaño fijo, operaciones para leer y escribir, elementos por índice etc \cite{tda}.

Todos estos conceptos serán útiles para facilitar el análisis y tener una buena implementación de estos para la solución de los problemas propuestos.

\section{Desarrollo}

Descripción de la implementación realizada y pruebas efectuadas.

\section{Resultados}

Presentación de capturas de pantalla y descripción de resultados.

\section{Conclusiones}

Análisis de los resultados obtenidos y la aplicación de los conceptos teóricos.

\begin{thebibliography}{9}

\bibitem{algoritmo}
R. Maluenda, ``Qué es un algoritmo informático: características, tipos y ejemplos,'' Profile, 2026. [En línea]. Disponible en: https://profile.es/blog/que-es-un-algoritmo-informatico/. [Accedido: 15-feb-2026].

\bibitem{arreglos}
ONMEX, ``¿Qué son los arreglos en programación?,'' 2026. [En línea]. Disponible en: https://onmex.mx/tecnologia-y-desarrollo/que-son-los-arreglos-en-programacion/. [Accedido: 15-feb-2026].

\bibitem{tda}
Wikipedia, ``Tipo de dato abstracto,'' 2026. [En línea]. Disponible en: https://es.wikipedia.org/wiki/Tipo_de_dato_abstracto. [Accedido: 15-feb-2026].

\end{thebibliography}

\end{document}
